\documentclass[12pt,a4paper]{article}
\usepackage[utf8]{inputenc}
\usepackage[indonesian]{babel}
\usepackage{amsmath}
\usepackage{amsfonts}
\usepackage{amssymb}
\usepackage{geometry}

\geometry{margin=2.5cm}

\title{Penyelesaian Luas Bidang Polar Non-Irisan}
\author{Ahmad Fakhrul Bawani}
\date{}

\begin{document}

\maketitle

\section*{Soal}
Diketahui dua persamaan kurva dalam koordinat polar:
\begin{align*}
  r_1 &= 2 - 2\sin\theta \quad \text{(Kardioida menghadap bawah)} \\
  r_2 &= 2 + 2\sin\theta \quad \text{(Kardioida menghadap atas)}
\end{align*}

Akan dicari luas daerah yang dibatasi oleh kedua kurva tersebut tetapi \textbf{bukan merupakan irisan} (symmetric difference) dari keduanya.

\subsection*{1. Rumus Utama Luas dalam Koordinat Polar}
Secara umum, rumus elemen luas daerah ($dA$) dalam koordinat polar yang disapu antara dua kurva $r_{\text{luar}}$ dan $r_{\text{dalam}}$ terhadap perubahan sudut $d\theta$ adalah:
\[
dA = \frac{1}{2} \left( r_{\text{luar}}^2 - r_{\text{dalam}}^2 \right) d\theta
\]
Sehingga, luas daerah non-irisan pada batas sudut $\alpha \le \theta \le \beta$ ditentukan oleh rumus integral:
\[
A = \int_{\alpha}^{\beta} \frac{1}{2} \left( r_{\text{luar}}^2 - r_{\text{dalam}}^2 \right) d\theta
\]

\subsection*{2. Menentukan Titik Potong Kedua Kurva (Batas Sudut)}
Untuk menentukan batas integral daerah, kita mencari titik potong dengan menyamakan kedua fungsi ($r_1 = r_2$):
\[
2 - 2\sin\theta = 2 + 2\sin\theta
\]
\[
-4\sin\theta = 0 \implies \sin\theta = 0
\]
Diperoleh nilai sudut pemotong pada satu putaran penuh adalah:
\[
\theta = 0 \quad \text{dan} \quad \theta = \pi
\]
Titik potong ini membagi wilayah grafik secara simetris tepat di sepanjang sumbu-x (garis polar horizontal).

\subsection*{3. Strategi Perhitungan Menggunakan Sifat Simetri}
Berdasarkan visualisasi grafis, wilayah non-irisan terdiri dari dua daun telinga yang simetris:
\begin{itemize}
  \item Daun bagian atas: dibatasi oleh kurva luar $r_2$ dan kurva dalam $r_1$ pada rentang $\theta \in [0, \pi]$.
  \item Daun bagian bawah: dibatasi oleh kurva luar $r_1$ dan kurva dalam $r_2$ pada rentang $\theta \in [\pi, 2\pi]$.
\end{itemize}
Karena kedua area tersebut memiliki ukuran geometri yang identik akibat simetri radial, luas total non-irisan dapat dihitung dengan mencari luas daun bagian atas ($A_{\text{atas}}$) kemudian dikalikan dua:
\[
A_{\text{total}} = 2 \times A_{\text{atas}}
\]

\subsection*{4. Menghitung Luas Daun Bagian Atas ($A_{\text{atas}}$)}
Pada rentang $0 \le \theta \le \pi$, kurva $r_2$ bertindak sebagai batas luar ($r_{\text{luar}}$) dan $r_1$ sebagai batas dalam ($r_{\text{dalam}}$):
\[
A_{\text{atas}} = \frac{1}{2} \int_{0}^{\pi} \left( r_2^2 - r_1^2 \right) d\theta
\]
Jabarkan pengurangan komponen kuadrat di dalam integral:
\[
r_2^2 - r_1^2 = (2 + 2\sin\theta)^2 - (2 - 2\sin\theta)^2
\]
\[
r_2^2 - r_1^2 = (4 + 8\sin\theta + 4\sin^2\theta) - (4 - 8\sin\theta + 4\sin^2\theta)
\]
Komponen konstanta dan $\sin^2\theta$ habis saling mengurangi, sehingga menyisakan:
\[
r_2^2 - r_1^2 = 16\sin\theta
\]
Substitusikan kembali hasil penyederhanaan ke dalam integral:
\[
A_{\text{atas}} = \frac{1}{2} \int_{0}^{\pi} 16\sin\theta \, d\theta = 8 \int_{0}^{\pi} \sin\theta \, d\theta
\]
Lakukan proses integrasi:
\[
A_{\text{atas}} = 8 \Big[ -\cos\theta \Big]_{0}^{\pi}
\]
Evaluasi nilai menggunakan batas atas ($\pi$) dan batas bawah ($0$):
\[
A_{\text{atas}} = 8 \Big( -\cos(\pi) - (-\cos(0)) \Big) = 8 \Big( 1 + 1 \Big) = 16
\]

\subsection*{5. Menghitung Luas Total Non-Irisan ($A_{\text{total}}$)}
Sesuai strategi simetri di awal, kalikan luas daun atas tersebut dengan dua untuk mendapatkan total seluruh daerah non-irisan:
\[
A_{\text{total}} = 2 \times A_{\text{atas}} = 2 \times 16 = 32
\]

\subsection*{Kesimpulan}
Luas daerah yang dibatasi oleh kurva $r = 2 - 2\sin\theta$ dan $r = 2 + 2\sin\theta$ yang bukan merupakan irisan keduanya adalah:
\[
A_{\text{total}} = 32 \quad \text{satuan luas}
\]

\end{document}